\Chapter{Results}
\label{Ch:Results}

The results chapter is for presenting and discussing your findings, which can split into sections if the experiment has multiple parts, or stages.  You might also want to have a `Discussion' or `Discussion of Results' chapter, which may focus on either more detailed discussion of particular results, or more comprehensive discussion of the results and system performance as a whole.  If you have a Discussion of Results section, then you do not want to discuss too much about the specific results in this chapter and rather move the main discussion or reflective considerations of these to the Discussion of Results chapter.

However, and this is an important reminder, ensure that you do have some text of discussion for your results, to take the reader through your results and the figures you may provide.  Make sure not to just put one figure after another without any attempt to explain the sequencing and what is being shown and perhaps some key issues in the figures presented (a monolithic sequence of figure after figure without any attempt at explanation will get undesirable responses from examiners).

If you have an experiments section, it is often useful to have a clear connection of experiment section corresponding to a result section showing results for that test. Examiners often appreciate that sort of clear and consistent structure that is easy to follow.  For example, if you have section~5.1 as Modular Testing, then you could correspond to~6.1 Modular Testing Results (the two sections should be cross-linked with \verb|\label| and \verb|\ref|, in both directions)

\section{Tips for Results Figures}

Include good quality graphs (see Fig.~\ref{fig:r_vs_N_f=0_0005_P=90} for an example).  These were produced by the Octave code presented in listings~\ref{lst:FormatFig} and~\ref{lst:PlotExample}.  You can play around with the \texttt{PaperSize} and \texttt{PaperPosition} variables to change the aspect ratio.  An easy way to obtain more space for an article is to use wide, flat figures, such as Fig.~\ref{fig:Line_Signals}.

\Figure[width=0.8\columnwidth]{The correlation coefficient as a function of sample count.}{r_vs_N_f=0_0005_P=90}

\begin{Matlab_float}{Octave function to format a figure and save it to a high quality PDF graph}{FormatFig}
  function FormatFig(X, Y, File);
    set(gcf, 'PaperUnits'      , 'inches');
    set(gcf, 'PaperOrientation', 'landscape');
    set(gcf, 'PaperSize'       ,       [8, 4]);
    set(gcf, 'PaperPosition'   , [0, 0, 8, 4]);
   
    set(gca, 'FontName', 'Times New Roman');
    set(gca, 'Position', [0.1 0.2 0.85 0.75]);
   
    xlabel(["\n" X]);
    ylabel([Y "\n\n"]);
   
    setenv("GSC", "GSC"); # Eliminates stupid warning
    print(...
      [File '.pdf'],...
      '-dpdf'...
    );
  end
\end{Matlab_float}

\begin{Matlab_float}{Example of how to use the FormatFig function}{PlotExample}
  figure;                                   # Create a new figure
  # Some code to calculate the various variables to plot...
  plot(N, r, 'k', 'linewidth', 4); grid on; # Plot the data
  xlim([0 360]);                            # Limit the x range
  ylim([-1 1]);                             # Limit the y range
  set(gca, 'xtick', [0 90 180 270 360]);    # Set the x labels

  FormatFig(...                             # Call the function with:
    'Phase shift [\circ]',...                      # The x title
    'Correlation coefficient',...                  # The y title
    ['r_vs_N;_f=' num2str(f) ';_P=' num2str(P)]... # Format the file name
  );
  close all;                                # Close all open figures
\end{Matlab_float}

\Figure[width=0.8\columnwidth]{Oscilloscope measurement showing physical line signals on both ends of a transmission line during master switch-over~\cite{Taylor_2016}.}{Line_Signals}

Always remember to include axes text, units and a meaningful caption in your graphs.  When typing units, a \micro{} sign has a tail!  The letter ``u'' is not a valid unit prefix.  When typing resistor values, use the \ohm{} symbol.

\section{Tables}

Tables are often a convenient means by which to specify lists of parameters.  An example table is presented in table~\ref{tab:Example}. You can use \href{https://www.tablesgenerator.com/}{Tablesgenerator} to make your \LaTeX\ tables.

\Table{My Informative Table}{lcr}{ % This format specifies 3 columns with left, centre and right allignment
  \hline
  \hline
  \textbf{Heading 1} & \textbf{Heading 2} & \textbf{Heading 3} \\
  \hline
  Data & 123 & 321 \\
  Data & 456 & 654 \\
  Data & 789 & 987 \\
  \hline
  \hline
}{Example}

\section{Pictures and Screen-shots}

When you include screen-shots, pdf\LaTeX{} supports JPG and PNG file formats.  PNG is preferred for screen-shots, as it is a loss-less format.  JPG is preferred for photos, as it results in a smaller file size.  It's generally a good idea to resize photos (not screen-shots) to be no more that 300~dpi, in order to reduce file size.  For 2-column article format papers, this translates to a maximum width of 1024.  \textbf{Never change the aspect ratio of screen-shots and pictures!}

It is highly recommended to make use of the \verb|\Figure| macro for figures.  It puts all the formatting tweaks in one place, so that you don't need to update all the individual figure inclusion points when you want to do a styling update.  The file name is used for the \LaTeX\ label, such as ``Fig.~\ref{fig:UCT.pdf}''.

\Figure[width=0.6\columnwidth]{An example image with custom scaling}{UCT.pdf}

Make sure to always use the best quality image possible.  Use JPEG for photos, PNG for screen-shots and PDF (scalable vector graphics) for everything else.  JPEG is lossy, but good for photos, whereas PNG is lossless and good for images with large areas of solid colour, as can be seen in Fig.~\ref{fig:FigureQuality}.

\begin{FigureEnvironment}{Comparison of various image format qualities}%
  \label{fig:FigureQuality}%
  \hfill
  \begin{minipage}[b]{0.3\textwidth}\centering\setlength{\parindent}{0mm}
    \includegraphics[width=\textwidth]{../Figures/UCT.jpg}\\%
    {\small (a)~~JPEG}%
  \end{minipage}
  \hfill
  \begin{minipage}[b]{0.3\textwidth}\centering\setlength{\parindent}{0mm}
    \includegraphics[width=\textwidth]{../Figures/UCT.png}\\%
    {\small (b)~~PNG}%
  \end{minipage}
  \hfill
  \begin{minipage}[b]{0.3\textwidth}\centering\setlength{\parindent}{0mm}
    \includegraphics[width=\textwidth]{../Figures/UCT.pdf}\\%
    {\small (c)~~SVG}%
  \end{minipage}
  \hfill\mbox{}\\[1mm]
\end{FigureEnvironment}%

\section{Maths}

\LaTeX{} has a very sophisticated maths rendering engine, as illustrated by equation~\ref{eq:Example}.  When talking about approximate answers, never use $\pm{54}$~V, as this implies ``positive or negative~54~V''.  Use $\approx{54}$~V or $\sim{54}$~V instead.

\begin{equation}
  y = \int_0^\infty e^{x^2} \mathrm{dx}
  \label{eq:Example}
\end{equation}


